\uuid{ieO8}
\titre{Intégration par parties}
\auteur{hanani ; mzari}
\organisation{exo7}
\licence{CC BY-NC-SA 4.0}
\datecreate{2020-01-20}
\niveau{L1}
\module{Analyse}
\chapitre{Calcul d'intégrales}
\sousChapitre{Primitives diverses}
\difficulte{3}
\theme{Primitives, intégration par parties}
\exo7qcmid{lille-2.primitives.037}

\contenu{
\begin{question}
\qclassification{127.11}
\qtags{motcle=intégration par parties}
Parmi les affirmations suivantes, cocher celles qui sont vraies :
\begin{answers}
\bad{$\displaystyle \int \frac{\mathrm{d}x}{\sqrt{1-x^2}}=-\arcsin (x)+k$, $k\in \Rr$.}
\good{$\displaystyle \int \sqrt{1-x^2}\mathrm{d}x=x\sqrt{1-x^2}+\int \frac{x^2\mathrm{d}x}{\sqrt{1-x^2}}$.}
\bad{L'égalité $x^2=(x^2-1)+1$ donne : $\displaystyle \int \frac{x^2\mathrm{d}x}{\sqrt{1-x^2}}=-\arcsin x-\int \sqrt{1-x^2}\mathrm{d}
x$.}
\good{Une primitive de $\displaystyle \sqrt{1-x^2}$ sur $]-1,1[$ est $\displaystyle \frac{1}{2}x\sqrt{1-x^2}+\frac{1}{2}\arcsin (x)$.}
\end{answers}
\vskip2mm
\begin{explanations}
D'abord $\displaystyle \int \frac{\mathrm{d}x}{\sqrt{1-x^2}}=\arcsin (x)+k$, $k\in \Rr$. Une intégration par parties avec $u=\sqrt{1-x^2}$ et $v=x$ donne
$$\int \sqrt{1-x^2}\mathrm{d}x=x\sqrt{1-x^2}-\int \frac{-x^2\mathrm{d}x}{\sqrt{1-x^2}}.$$
Ensuite, en écrivant $\displaystyle -x^2=(1-x^2)-1$, on obtient :
$$\int \sqrt{1-x^2}\mathrm{d}x=x\sqrt{1-x^2}-\int \sqrt{1-x^2}\mathrm{d}x+\int \frac{\mathrm{d}x}{\sqrt{1-x^2}}.$$
D'où : $\displaystyle 2\int \sqrt{1-x^2}\mathrm{d}x=x\sqrt{1-x^2}+\arcsin x+k$.
\end{explanations}
\end{question}
}
